\begin{summarybox}
\textbf{一张速记卡: 一式、两轴、两种证明}

\textbf{核心一式:}
\[
\boxed{\sqrt{r_1r_2}\,d=ab},\qquad
\boxed{d=\frac{ab}{\sqrt{a^2-e^2x_0^2}}}.
\]
$r_1,r_2$ 是\textbf{切点到焦点}的距离;
$d$ 是\textbf{中心到切线}的垂直距离. 两种距离千万别混.

\textbf{最值两轴:}
\[
\begin{array}{c|c|c}
\text{切点} & r_1r_2 & d\\ \hline
(\pm a,0)\ \text{长轴顶点} & b^2 & \boxed{a}\ \text{最大}\\
(0,\pm b)\ \text{短轴顶点} & a^2 & \boxed{b}\ \text{最小}
\end{array}
\]
故 $b\le d\le a$.

\textbf{最快代数法 (3 个动作):}
\[
\begin{gathered}
l:\ \frac{x_0x}{a^2}+\frac{y_0y}{b^2}=1,\qquad
\frac1{d^2}=\frac{x_0^2}{a^4}+\frac{y_0^2}{b^4},\\
\frac{a^2b^2}{d^2}
=\frac{b^2x_0^2}{a^2}+\frac{a^2y_0^2}{b^2}
=a^2-\frac{c^2x_0^2}{a^2}=r_1r_2,\\
\therefore\quad \sqrt{r_1r_2}\,d=ab.
\end{gathered}
\]
记忆口诀: \textbf{写切线、取距离、代椭圆、认乘积}.

\textbf{最快几何联想:} 设 $\varphi=\angle F_1PF_2$,
\[
\boxed{d=a\cos\frac\varphi2},\qquad
\boxed{\sqrt{r_1r_2}\cos\frac\varphi2=b}.
\]
两式相除便得主公式. 依据是\textbf{焦点三角形 + 切线等角反射 + 余弦定理}.

\textbf{与圆的对应:} 圆有 $d=R$, 椭圆有
$\sqrt{r_1r_2}\,d=ab$.
令 $a=b=R$、两焦点重合, 就得到 $R\cdot R=R^2$.

\textbf{姊妹性质:} 若 $h_1,h_2$ 是\textbf{两焦点到切线}的垂距,
则 $h_1h_2=b^2$ 且 $d=(h_1+h_2)/2$.

\textbf{高考识别}: ``椭圆 + 切线 + 中心距 + 焦半径'' 直接用 C012;
给 $r_1-r_2$ 时先用
$4r_1r_2=(r_1+r_2)^2-(r_1-r_2)^2$.
\end{summarybox}
